\subsection{Ellipse from the polar formula}

Assume
\[
0<e<1.
\]
Then
\[
1-e\leq1+e\cos\theta\leq1+e,
\]
so the denominator is positive for every angle. The radius is finite in every
direction and the curve closes around the focus.

The nearest point occurs at $\theta=0$:
\[
\boxed{r_{\min}=\frac{p}{1+e}}.
\]
The farthest point occurs at $\theta=\pi$:
\[
\boxed{r_{\max}=\frac{p}{1-e}}.
\]
Thus increasing $e$ makes the two extreme focal distances more unequal.

\begin{center}
\begin{tikzpicture}[scale=0.82,>=stealth]
    \draw[axis,->] (-6.0,0)--(5.3,0) node[right] {polar axis};
    \draw[axis,->] (0,-3.5)--(0,3.7);
    \draw[blue!70!black,line width=1.1pt]
        plot[domain=0:360,samples=220]
        ({2.4*cos(\x)/(1+0.25*cos(\x))},{2.4*sin(\x)/(1+0.25*cos(\x))});
    \draw[orange!85!black,line width=1.1pt]
        plot[domain=0:360,samples=220]
        ({2.4*cos(\x)/(1+0.65*cos(\x))},{2.4*sin(\x)/(1+0.65*cos(\x))});
    \fill[geometricSubsetColor] (0,0) circle (2.2pt) node[below left] {focus};
    \node[blue!70!black] at (-3.0,2.6) {$e=0.25$};
    \node[orange!85!black] at (-5.0,-2.6) {$e=0.65$};
\end{tikzpicture}
\end{center}

For a Cartesian ellipse with semimajor axis $a$ and semiminor axis $b$,
\[
e=\sqrt{1-\frac{b^2}{a^2}},
\qquad
p=a(1-e^2)=\frac{b^2}{a}.
\]
Therefore the focus-centred form is
\[
\boxed{r(\theta)=\frac{a(1-e^2)}{1+e\cos\theta}}.
\]

\begin{interpretationbox}[Why the focus-centred form matters]
The Cartesian equation displays the centre and semiaxes. The polar equation
displays the distance from a focus as the direction changes. This second view is
the one that later matches central-force motion.
\end{interpretationbox}
